% Source: https://artofproblemsolving.com/community/c3h1708605_tree_planting_permutation
% Stable identity: aops-topic-1708605-q01
Place the $10$ maples first; they create $11$ gaps (including the two ends). The oaks and birches go into the gaps, and inside a gap no oak may touch a birch, so each gap receives trees of only one kind (or none).

Suppose the oaks occupy $i$ gaps and the birches $j$ gaps ($1\le i,j\le3$). Splitting $3$ identical oaks into $i$ nonempty groups in order can be done in $\dbinom{2}{i-1}$ ways, and the same for birches. Choosing the gaps: $\dbinom{11}{i}\dbinom{11-i}{j}$ ways. So the count is
\[
\sum_{i=1}^{3}\sum_{j=1}^{3}\binom{2}{i-1}\binom{2}{j-1}\binom{11}{i}\binom{11-i}{j}.
\]
Row by row:
\begin{itemize}
\item $i=1$: $11\,(10+2\cdot45+120)=11\cdot220=2420$;
\item $i=2$: $2\cdot55\,(9+2\cdot36+84)=110\cdot165=18150$;
\item $i=3$: $165\,(8+2\cdot28+56)=165\cdot120=19800$.
\end{itemize}
The total is $2420+18150+19800=\boxed{40370}$.
